% Part: first-order-logic
% Chapter: natural-deduction
% Section: provability-quantifiers

% verification of properties of provability needed for maximally
% consistent sets in the completeness chapter.

\documentclass[../../../include/open-logic-section]{subfiles}

\begin{document}

\olfileid[hi]{fol}{ntd}{qpr}

\olsection{\usetoken{S}{derivability} और परिमाणक}

\begin{explain}
  पूर्णता प्रमेय के लिए यह भी आवश्यक है कि स्वाभाविक निगमन के नियम इस अनुभाग
  में स्थापित~$\Proves$ से संबंधित तथ्य प्रदान करें।
\end{explain}

\begin{thm}
\ollabel{thm:strong-generalization} यदि $c$ ऐसा अचर है जो $\Gamma$ या
$!A(x)$ में नहीं आता, और $\Gamma \Proves !A(c)$, तो $\Gamma
\Proves \lforall[x][!A(x)]$।
\end{thm}

\begin{proof}
$\delta$ को $!A(c)$ की $\Gamma$ से !!a{derivation} मानें।
\Intro{\lforall} अनुमान जोड़ने पर $\lforall[x][!A(x)]$ की !!a{derivation}
मिलती है। चूँकि $c$, $\Gamma$ या $!A(x)$ में नहीं आता, इसलिए
अभिलाक्षणिक-चर शर्त पूरी होती है।
\end{proof}

\begin{prop}
\ollabel{prop:provability-quantifiers}
\begin{tagenumerate}{prvEx,prvAll}
\tagitem{prvEx}{$!A(t) \Proves \lexists[x][!A(x)]$.}{}

\tagitem{prvAll}{$\lforall[x][!A(x)] \Proves !A(t)$.}{}
\end{tagenumerate}
\end{prop}

\begin{proof}
\begin{tagenumerate}{prvEx,prvAll}
\tagitem{prvEx}{निम्नलिखित !!a{derivation} है—
  ~$\lexists[x][!A(x)]$ की~$!A(t)$ से:
\begin{prooftree}
\AxiomC{$!A(t)$}
\RightLabel{\Intro{\lexists}}
\UnaryInfC{$\lexists[x][!A(x)]$}
\end{prooftree}}{}

\tagitem{prvAll}{निम्नलिखित !!a{derivation} है—~$!A(t)$ की
  ~$\lforall[x][!A(x)]$ से:
\begin{prooftree}
\AxiomC{$\lforall[x][!A(x)]$}
\RightLabel{\Elim{\lforall}}
\UnaryInfC{$!A(t)$}
\end{prooftree}}{}
\end{tagenumerate}
\end{proof}

\end{document}
